%============================================================ 1. 서론 (압축판)
\section{서론}
\label{sec:intro}

\begin{figure*}[tp]
\centering
\includegraphics[width=0.72\textwidth]{fig4.png}
\caption{제안 체계의 2계층 구조}
\label{fig:architecture}
\end{figure*}

다수의 저비용 무인수상정(USV)이 고가치 함정을 향해 여러 방향에서 동시에 돌진하는 자폭
공격은 비대칭 위협의 전형이다. 유한한 방어 자산의 배분은 무기-표적 할당 문제로 오래
연구되었으나~\citep{ahuja2007wta}, 표적 수가 요격 자산을 넘어서면 최적 할당으로도 남는
표적이 생긴다. 교란에 의한 연성 대응(soft-kill)은 위협을 제거하지 못하고 유선 조종 표적에
무력하며~\citep{wang2021cuas, sutea2026fiber},
직접 타격하는 경성 대응(hard-kill)은 값싼 표적에 고가 유도탄을 쓰는 비용
구조~\citep{rumbaugh2024cost}와 표적 수에 따른 포화를 피하지 못한다.

본 연구는 저비용 그물에 의한 물리적 포획을 택한다~\citep{wang2021cuas}. 이 방식은 통신
방식과 무관하게 표적을 정지시킨다. 방어정은 적선보다 느려 추격할 수 없으므로, 적선이
지나갈 접근 회랑을 앞질러 그물벽으로 차단한다.

이 문제는 성격이 다른 두 판단을 요구한다. 하나는 어느 방어정이 어느 적 무리를 맡는가라는
이산적 배정이고, 다른 하나는 맡은 무리를 해면 어디에서 막는가라는 연속적 기동이다.
둘은 같은 주기로 갱신되지만 산출에 걸리는 시간이 크게 다르다. 본 연구는 느린 배정을
언어모델 지휘관에게, 빠른 기동을 학습된 정책에 맡기고 느린 계층만 비동기로
분리한다(Fig.~\ref{fig:architecture}).

학습과 평가는 OpenStreetMap 해안선을 래스터화한 실해역 지형에서 수행한다. 관측은 모선 중심
격자의 다채널 영상이고, 행동은 그 격자의 유효 픽셀 하나를 고르는 것이라 불가능한 자리가 애초에
뽑히지 않는다. 출력은 제어 신호가 아니라 경유점이라 저수준 추종과 분리된다(\ref{sec:method}절).
각 계층이 무엇을 얼마나 개선하는지는 시드 대응 $2\times2$ 설계로 분리해 측정한다
(\ref{sec:experiment}절).
